\documentclass[11pt]{article}
\usepackage[utf8]{inputenc}
\usepackage[spanish]{babel}
\usepackage{amsmath,amssymb,amsthm}
\usepackage{graphicx}
\usepackage[margin=2.5cm]{geometry}
\usepackage{enumitem}

\newtheorem{definicion}{Definición}[section]
\newtheorem{teorema}{Teorema}[section]
\newtheorem{proposicion}{Proposición}[section]
\newtheorem{corolario}{Corolario}[section]
\newtheorem{axioma}{Axioma}[section]
\theoremstyle{remark}
\newtheorem*{ejemplo}{Ejemplo}
\newtheorem*{observacion}{Observación}

\title{Semana 2 -- Materia}
\author{Tutorías -- Introducción al Cálculo}
\date{10/08--14/08/2026}

\begin{document}
\maketitle

\section{Teorema del factor y aplicaciones}

Recordemos de la tutoría anterior que, para un polinomio $P$ y un número
$\alpha$, se tiene que $\alpha$ es raíz de $P$ si y solamente si $P$ es
divisible por $(x-\alpha)$. Esta equivalencia (a veces llamada
\textbf{Teorema del Factor}) tiene consecuencias importantes sobre cuántas
raíces puede tener un polinomio y sobre cómo encontrarlas.

\subsection{Cantidad de raíces de un polinomio}

\begin{teorema}\label{teo:raices}
Si un polinomio $P(x) = a_0 + a_1x + a_2x^2 + \cdots + a_nx^n$, con $a_n \neq 0$,
tiene más de $n$ raíces distintas, entonces necesariamente todos sus
coeficientes son nulos (es decir, $P(x)=0$ para todo $x$).
\end{teorema}

La idea de la demostración es aplicar repetidamente la relación raíz-divisibilidad:
si $\alpha_1$ es raíz de $P$, entonces $P(x) = (x-\alpha_1)Q(x)$ con $Q$ de
grado menor. Si además $\alpha_2 \neq \alpha_1$ es raíz de $P$, entonces
$P(\alpha_2) = (\alpha_2-\alpha_1)Q(\alpha_2) = 0$, y como $\alpha_2 \neq
\alpha_1$, debe ser $\alpha_2$ raíz de $Q$. Repitiendo este argumento con
$\alpha_1,\dots,\alpha_n$ raíces distintas de $P$, se llega a
$$P(x) = (x-\alpha_1)(x-\alpha_2)\cdots(x-\alpha_n)\,k$$
para alguna constante $k$. Si hubiera todavía otra raíz distinta de las
anteriores, forzosamente $k=0$, y al desarrollar el producto todos los
coeficientes de $P$ resultan nulos. En otras palabras: \textbf{un polinomio no
nulo de grado $n$ tiene a lo más $n$ raíces distintas}.

\begin{corolario}
Dos polinomios toman los mismos valores para cada $x$ si y solamente si sus
coeficientes homólogos son iguales.
\end{corolario}

En efecto, si dos polinomios $P$ y $Q$ coinciden para todo $x$, entonces el
polinomio diferencia $P-Q$ tiene infinitas raíces (todos los reales), muchas
más que su grado, así que por el Teorema~\ref{teo:raices} todos sus
coeficientes deben ser $0$, es decir, los coeficientes de $P$ y $Q$ coinciden.
El recíproco es inmediato.

\begin{ejemplo}
Encontremos constantes $a, b \in \mathbb{R}$ tales que
$$x^3 + ax + b = (x-1)(x-2)(x+3) \quad \text{para todo } x.$$
Desarrollamos el lado derecho:
$$(x-1)(x-2) = x^2 - 3x + 2, \qquad (x^2-3x+2)(x+3) = x^3 - 7x + 6.$$
Como ambos polinomios deben coincidir para todo $x$, el Corolario anterior
nos dice que sus coeficientes homólogos deben ser iguales. Comparando
$x^3 + 0\cdot x^2 + ax + b$ con $x^3 + 0 \cdot x^2 - 7x + 6$, obtenemos
$a = -7$ y $b = 6$.
\end{ejemplo}

\subsection{Raíces racionales de un polinomio}

Cuando un polinomio tiene coeficientes enteros, el siguiente resultado acota
mucho la búsqueda de sus raíces racionales.

\begin{proposicion}\label{prop:raicesracionales}
Sea $P(x) = a_0 + a_1x + \cdots + a_nx^n$ un polinomio con coeficientes
enteros. Si $P$ tiene una raíz racional $\frac{p}{q}$, con $p,q$ enteros
primos relativos (sin factores comunes salvo $1$), entonces $p$ divide a
$a_0$ y $q$ divide a $a_n$.
\end{proposicion}

La idea es reemplazar $x=\frac{p}{q}$ en $P(x)=0$ y multiplicar por $q^n$
para obtener
$$a_0q^n + a_1pq^{n-1} + \cdots + a_np^n = 0.$$
Despejando el término $a_np^n$ al lado derecho, el resto de la suma es
múltiplo de $q$, así que $q$ divide a $a_np^n$; como $p$ y $q$ son primos
relativos, $q$ debe dividir a $a_n$. De forma simétrica, despejando el
término $a_0q^n$, se obtiene que $p$ divide a $a_0$.

Esta proposición \textbf{no garantiza} que esos candidatos $\frac{p}{q}$ sean
efectivamente raíces: solo reduce la búsqueda a una lista finita de
candidatos, que luego hay que verificar uno por uno evaluando el polinomio
(o usando el Teorema del Resto).

\begin{ejemplo}
Busquemos las raíces racionales de $P(x) = 2x^3 - x^2 - 13x - 6$.

Los divisores de $a_0 = -6$ son $\pm1,\pm2,\pm3,\pm6$, y los divisores de
$a_3 = 2$ son $\pm1,\pm2$. Los candidatos $\frac{p}{q}$ son entonces
$$\pm1,\ \pm2,\ \pm3,\ \pm6,\ \pm\tfrac12,\ \pm\tfrac32.$$
Probamos $x=3$: $P(3) = 2(27) - 9 - 39 - 6 = 54-9-39-6=0$. Así, $3$ es raíz.
Dividiendo $P(x)$ por $(x-3)$ con coeficientes indeterminados ($\alpha=3$,
$a_3=2,a_2=-1,a_1=-13,a_0=-6$):
$$c_3=2,\quad c_2=-1+3(2)=5,\quad c_1=-13+3(5)=2,\quad b_0=-6+3(2)=0,$$
obtenemos $Q(x) = 2x^2+5x+2$. Factorizando (o usando la fórmula
cuadrática), $2x^2+5x+2 = (2x+1)(x+2)$, con raíces $x=-\tfrac12$ y $x=-2$.
Por lo tanto,
$$P(x) = (x-3)(2x+1)(x+2),$$
y las raíces racionales de $P$ son $3,\ -\tfrac12,\ -2$ — todas efectivamente
dentro de la lista de candidatos que predecía la Proposición~\ref{prop:raicesracionales}.
\end{ejemplo}

\begin{observacion}
El libro menciona, sin demostración (por ser demasiado avanzada para este
curso), un resultado general conocido como el \textbf{Teorema Fundamental
del Álgebra}: todo polinomio no constante tiene al menos una raíz, si se
permite ampliar el conjunto de números donde se buscan las raíces más allá
de los reales. En este curso no hemos estudiado ese conjunto ampliado de
números, así que por ahora basta con tener presente que \emph{no todo
polinomio tiene raíces reales} (por ejemplo, $P(x) = x^2+1$ no se anula para
ningún $x$ real), y que las herramientas de esta sección (Teorema del
Factor, Proposición~\ref{prop:raicesracionales}) sirven para encontrar las
raíces \emph{reales} (en particular las racionales) cuando existen.
\end{observacion}

\section{Axiomas de Orden I}

Así como los números reales satisfacen axiomas para la suma y el producto,
también satisfacen axiomas que describen la relación "menor que" ($<$). A
partir de estos axiomas se puede demostrar cualquier otra propiedad del
orden.

\subsection{Los axiomas}

\begin{axioma}[Orden lineal]
Para todo $x,y \in \mathbb{R}$, si $x \neq y$ entonces $x<y$ o $y<x$.
\end{axioma}

\begin{axioma}[Asimetría]
Para todo $x,y \in \mathbb{R}$, si $x<y$ entonces no se tiene $y<x$.
\end{axioma}

\begin{axioma}[Transitividad]
Para todo $x,y,z \in \mathbb{R}$, si $x<y$ e $y<z$ entonces $x<z$.
\end{axioma}

\begin{axioma}[Compatibilidad con la suma]
Para todo $x,y,z \in \mathbb{R}$, si $x<y$ entonces $x+z<y+z$.
\end{axioma}

\begin{axioma}[Compatibilidad con el producto por positivos]
Para todo $x,y,z \in \mathbb{R}$, si $x<y$ y $0<z$ entonces $x\cdot z < y \cdot z$.
\end{axioma}

A partir de estos cinco axiomas se definen las demás relaciones de orden
habituales.

\begin{definicion}[Relaciones de orden]
Sean $x,y \in \mathbb{R}$. Entonces:
\begin{enumerate}[label=(\Roman*)]
\item $x>y \iff y<x$ \quad ($x$ es mayor que $y$).
\item $x\geq y \iff (x>y \ \text{o} \ x=y)$ \quad ($x$ es mayor o igual que $y$).
\item $x\leq y \iff (x<y \ \text{o} \ x=y)$ \quad ($x$ es menor o igual que $y$).
\item $x$ es \textbf{positivo} $\iff x>0$.
\item $x$ es \textbf{negativo} $\iff x<0$.
\end{enumerate}
\end{definicion}

\subsection{Consecuencias de los axiomas}

Con solo estos axiomas y definiciones se pueden demostrar muchas propiedades
que usamos constantemente al trabajar con desigualdades, entre ellas:

\begin{teorema}\label{teo:orden}
Para $x,y,z \in \mathbb{R}$:
\begin{enumerate}[label=(\Roman*)]
\item $\neg(x<x)$.
\item $\neg(x<y) \iff x \geq y$.
\item $x>0 \iff -x<0$.
\item $(x>0 \ \text{y} \ y>0) \Rightarrow x+y>0$.
\item $x>y \iff x-y>0$.
\item $(x>0 \ \text{y} \ y>0) \Rightarrow xy>0$.
\item $1>0$.
\item $x^2 \geq 0$.
\end{enumerate}
\end{teorema}

Veamos cómo se demuestran algunas de estas propiedades directamente a partir
de los axiomas, a modo de ejemplo.

\begin{ejemplo}[Demostración de (I)]
Supongamos, por contradicción, que $x<x$. Aplicando el axioma de asimetría
con $y=x$: si $x<x$ entonces no puede ser $x<x$, lo cual es una
contradicción inmediata. Luego $\neg(x<x)$.
\end{ejemplo}

\begin{ejemplo}[Demostración de (V)]
Queremos ver que $x>y \iff x-y>0$, es decir (usando la definición de $>$),
que $y<x \iff 0<x-y$.
\begin{itemize}
\item[$(\Rightarrow)$] Si $y<x$, sumamos $-y$ a ambos lados usando el axioma
de compatibilidad con la suma (con $z=-y$): $y+(-y) < x+(-y)$, es decir,
$0 < x-y$.
\item[$(\Leftarrow)$] Si $0<x-y$, sumamos $y$ a ambos lados: $0+y < (x-y)+y$,
es decir, $y<x$.
\end{itemize}
\end{ejemplo}

Este tipo de argumento (sumar o multiplicar por algo conveniente, usando
solo los axiomas) permite derivar propiedades nuevas. Por ejemplo:

\begin{ejemplo}
Demostremos que si $x<y$ entonces $-y<-x$. Partimos de $x<y$ y sumamos
$(-x-y)$ a ambos lados (axioma de compatibilidad con la suma, con
$z=-x-y$):
$$x+(-x-y) < y+(-x-y) \ \Longrightarrow\ -y < -x.$$
\end{ejemplo}

\begin{ejemplo}
Como ilustración numérica del axioma de compatibilidad con el producto:
sabemos que $2<5$ y que $3>0$; el axioma nos garantiza entonces que
$2\cdot 3 < 5 \cdot 3$, es decir, $6<15$, lo cual es evidentemente cierto.
\end{ejemplo}

\end{document}
